\PassOptionsToPackage{quiet}{fontspec}
\documentclass[a4paper,12pt,CJK,openany]{ctexbook}
\usepackage{amsfonts}
\usepackage{latexsym}
\usepackage{amsmath}
\usepackage{titlesec,titletoc,amsmath}
\usepackage{amssymb,amsmath,amsfonts,epsfig,fancyhdr,mathrsfs,amsthm,url,latexsym,dsfont,extarrows}
\usepackage{graphicx}
\usepackage[figuresright]{rotating}
% \usepackage{epstopdf}
\usepackage{booktabs}
\usepackage[nottoc]{tocbibind}%取消参考文献空格
\usepackage{geometry}
\headsep -0.8 cm
\geometry{left=3.0cm,right=3.0cm,top=2.0cm,bottom=2.5cm}

\usepackage{setspace}
\usepackage{float}

%%%%%%%%%%%%%%%%%%%%%%%%%%%%%%%
\usepackage{xeCJK} % xeCJK 默认调用 Fandol 系列字体
\usepackage[style = gb7714-2015]{biblatex}
\addbibresource{ref.bib}
\usepackage{subfig}
\usepackage[ruled,linesnumbered]{algorithm2e}
\usepackage{threeparttable}
\usepackage{multirow}
\pagestyle{fancy}
% \usepackage{hyperref}
% \hypersetup{
%     colorlinks=true, linkcolor=black, filecolor=black,
%     urlcolor=black, citecolor=green}

% \renewcommand{\algorithmicrequire}{\textbf{输入:}}
% \renewcommand{\algorithmicensure}{\textbf{输出:}}
% \setcounter{tocdepth}{1}
%\setlength{\headheight}{15pt}

\renewcommand{\theequation}{\arabic{chapter}.\arabic{section}.\arabic{equation}}
\newtheorem{definition}  {{\heiti 定义}}[section]
\newtheorem{lemma}  {{\heiti 引理}}[section]
\newtheorem{theorem} {{\heiti 定理}}[section]
\newtheorem{example}{{\heiti 例  }}[section]
\newtheorem{remark}  {{\heiti 注  }}[section]
\newtheorem{cor}{{\heiti 推论}}[section]
\newtheorem{proposition}{{\heiti 命题}}[section]
\newtheorem{alg}{{\heiti 算法}}[section]

\newcommand{\diag}{{\rm diag}}
\newcommand{\RE}{{\rm Re}}

\newcommand{\chuhao}{\fontsize{55pt}{\baselineskip}\selectfont}     % 字号设置
\newcommand{\xiaochuhao}{\fontsize{36pt}{\baselineskip}\selectfont} % 字号设置
\newcommand{\yihao}{\fontsize{28pt}{\baselineskip}\selectfont}      % 字号设置
\newcommand{\erhao}{\fontsize{21pt}{\baselineskip}\selectfont}      % 字号设置
\newcommand{\xiaoerhao}{\fontsize{18pt}{\baselineskip}\selectfont}  % 字号设置
\newcommand{\sanhao}{\fontsize{15.75pt}{\baselineskip}\selectfont}
\newcommand{\xiaosanhao}{\fontsize{15pt}{\baselineskip}\selectfont} % 字号设置
\newcommand{\sihao}{\fontsize{13pt}{\baselineskip}\selectfont}      % 字号设置
\newcommand{\xiaosihao}{\fontsize{12pt}{\baselineskip}\selectfont}  % 字号设置
\newcommand{\wuhao}{\fontsize{10.5pt}{\baselineskip}\selectfont}    % 字号设置
\newcommand{\xiaowuhao}{\fontsize{9pt}{\baselineskip}\selectfont}   % 字号设置
\newcommand{\liuhao}{\fontsize{7.875pt}{\baselineskip}\selectfont}  % 字号设置
\newcommand{\qihao}{\fontsize{5.25pt}{\baselineskip}\selectfont}    % 字号设置

% \usepackage[numbers,sort,compress]{natbib}
% \newcommand{\upcite}[1]{\textsuperscript{\textsuperscript{\cite{#1}}}}

\usepackage{amsthm,amsmath,amssymb,fancyhdr,color,graphicx}
\newcommand{\hh}{\hspace*{24 pt}}
\newcommand{\h}{\hspace*{6 pt}}
\newcommand{\ho}{\hspace*{3 pt}}
\newcommand{\hi}{\hspace*{4 pt}}
\newcommand{\hhh}{\hspace*{96 pt}}
\newcounter{rom}
\newcommand{\rR}{I \hspace*{-1.2mm} R}
\newcommand{\xd}[1]{\setcounter{rom}{#1}{\rm(\roman{rom})}}
\newcommand{\Xd}[1]{\setcounter{rom}{#1}{\rm\Roman{rom}}}
\titlespacing{\chapter}{0pt}{-10pt}{-10pt}
\titlespacing{\section}{0pt}{*1}{*1}
\titlespacing{\subsection}{0pt}{*1}{*1}
\titleformat{\chapter}[hang]{\heiti\centering\sanhao\bfseries}{第\thechapter 章}{0em}{}
\titleformat{\section}[hang]{\heiti\xiaosanhao\bfseries}{~\thesection~}{0em}{}
\titleformat{\subsection}[hang]{\heiti\sihao\bfseries}{~\thesubsection~}{0.5em}{}

\def\beq{\begin{equation}}
	\def\eeq{\end{equation}}
\makeatletter \@addtoreset{equation}{section} \makeatother
\allowdisplaybreaks[4]
\pagestyle{empty} %取消页码

%%%%%%%%%%%%%%%%%%%%%%%%%%%%%%%%%%%%%%%%%%%%%%%%%%%%%%%%%%%%%%%%%%%%%%%%%

\begin{document}
	
%%%%%%%%%%%%%%%%%%%%%%%%%%%论文题目%%%%%%%%%%%%%%%%%%%%%%%%%%%
\thispagestyle{empty}
%\topskip1cm \rightline{10126$-$30408017}
\quad\\
\xiaosihao 分类号\underline{ \quad F224.32~~  }\hspace{7.5cm}单位代码\underline{ \quad \quad 10618\quad\quad \,}   \\
密\quad 级\underline{ \quad  公开 \quad~~ }\hspace{7.45cm} 学\quad\quad 号\underline{ ~666666666666~ }
%\hh\ho
%D\ C \underline{\hspace{3cm}
    %}\hspace{7.5cm}编号\underline{\hspace{3cm}       }
\vskip 1.0cm

%\begin{minipage}{0.8\textwidth}
%\begin{flushright}
%\begin{figure}[H]
\begin{center}
    \includegraphics[width=0.7\textwidth]{Figure_0}\\
\end{center}
%\end{figure}
%\end{flushright}
%\end{minipage}

\centerline{\heiti \chuhao 硕士学位论文 } \vskip 2 cm
\begin{center}
    \songti \yihao 论文标题
\end{center}
\vskip 3.0cm
{\heiti\xiaosanhao
    \[\begin{array}{rl}
        \mbox{研~\,究~\,生~\,姓~\,\,名：}&\mbox{\underline{\quad \quad *** \quad \quad}}\\
        \mbox{导师姓名及职称：}&\mbox{\underline{\quad ***教授 \quad}}\\
        % &\mbox{ \quad 副教授}\\
    \end{array}\]
}

{\songti\xiaosihao
    \[\begin{array}{rl}
    \mbox{ 申请学位类别~} \mbox{\underline{\quad \, \heiti{理学硕士} \quad }}  
    \mbox{ 学位授予单位  } \mbox{\underline{~\,\heiti{~~~重庆交通大学~}~~\,}} \\
    \mbox{ 一级学科名称~} \mbox{\underline{\quad \,  \heiti{系统科学}  \quad }}  
    \mbox{论文提交日期} \mbox{\underline{\heiti{2024 年 03 月  31 日\,}}}  \\
    \mbox{二级学科名称~} \mbox{\underline{\quad \,  \heiti{********}  \quad }} 
    \mbox{论文答辩日期} \mbox{\underline{\heiti{2024 年 06 月  06 日\,}}}  \\
     \end{array}\]
}
 \centerline{\sanhao 2024 年 06 月}

\newpage
\thispagestyle{empty} %\fancyhead[OC]{}
%\begin{center}
% \xiaoerhao
%A Thesis Submitted to Chongqing Jiaotong University in
%Partial Fulfillment of the Requirements for the Degree of Master
% \end{center}
\topskip 70pt
\vskip 1cm

\indent
\vspace{0.7cm}
\begin{center}
    \bf\erhao English Title
\end{center}

\vskip 1cm

\begin{center}
    \sihao A Dissertation Submitted for the Master’s Degree of Science
\end{center}

\vskip 3cm
\begin{eqnarray*}
	% \nonumber to remove numbering (before each equation)
	& &\text{\bf\xiaosanhao Candidate：***} \\
	\quad\\
	& &\text{\bf\xiaosanhao Supervisor：Prof. ***} \\
	% & &\text{\sanhao ~~~~~~~~~~~~~~~~Associate Professor Zhao Kequan } \\
	% & &\text{\sanhao Major：Technical Economic and Management} \\
	% & &\text{\sanhao College：College of Management}
\end{eqnarray*}
\vskip 4cm

\begin{center}
    \xiaosanhao  Chongqing Jiaotong University, Chongqing, China
\end{center}

\newpage
\thispagestyle{empty} %\fancyhead[OC]{}
\topskip 30pt

\begin{center}
    \kaishu \sanhao  重庆交通大学学位论文原创性声明
\end{center}
\vspace{0.3cm}

本人郑重声明：所呈交的学位论文，是本人在导师的指导下，独立进行研究工作所取得的成果。除文中已经注明引用的内容外，本论文不包含任何其他个人或集体已经发表或撰写过的作品成果。对本文的研究做出重要贡献的个人和集体，均已在文中以明确方式标明。本人完全意识到本声明的法律结果由本人承担。
\vspace{0.3cm}

学位论文作者签名：     ~~~~~~~~~~~~~~~~~~~~~~~~~~~~~~~~~~~~~~~~~~~~~~            日期： ~~~    年  ~~ 月 ~~~   日

\vspace{1cm}
\begin{center}
    \kaishu \sanhao  重庆交通大学学位论文版权使用授权书
\end{center}
\vspace{0.3cm}

本学位论文作者完全了解学校有关保留、使用学位论文的规定，同意学校保留并向国家有关部门或机构送交论文的复印件和电子版，允许论文被查阅和借阅。本人授权重庆交通大学可以将本学位论文的全部内容编入有关数据库进行检索，可以采用影印、缩印或扫描等复制手段保存和汇编本学位论文。同时授权中国科学技术信息研究所将本人学位论文收录到《中国学位论文全文数据库》，并进行信息服务（包括但不限于汇编、复制、发行、信息网络传播等），同时本人保留在其他媒体发表论文的权利。

\vspace{0.3cm}


学位论文作者签名：    ~~~~~~~~~~~~~~~~~~~~~~~~~~~~~~~~~~~~~~~~~~~~~~~~~           指导教师签名：~~~~~~~~~~~

\,日期：  ~~~ 年 ~~   月  ~~~  日  ~~~~~~~~~~~~~~~~~~~~~~~~~~~~~~~~~~~~~~~~~       日期： ~~~    年  ~~ 月 ~~~   日

{..................................................................................................................}

\vspace{0.3cm}
本人同意将本学位论文提交至中国学术期刊（光盘版）电子杂志社CNKI 系列数据库中全文发布，并按《中国优秀博硕士学位论文全文数据库出版章程》规定享受相关权益。

\vspace{0.3cm}

学位论文作者签名：   ~~~~~~~~~~~~~~~~~~~~~~~~~~~~~~~~~~~~~~~~~~~~~~~             指导教师签名：~~~~~~~~~~~

\,日期：  ~~~ 年 ~~   月  ~~~  日  ~~~~~~~~~~~~~~~~~~~~~~~~~~~~~~~~~~~~~~  ~      日期： ~~~    年  ~~ 月 ~~~   日




\setcounter{page}{0}
\newpage

%%%%%%%%%%%%%%%%%%%%%%%%%%%%%
\pagestyle{plain}
%\sloppy \baselineskip 20pt \pagestyle{fancy}
\fancyhead{}% clear all fields
\topskip 54pt
%\thispagestyle{fancy} %\fancyhead[CE]{\songti \wuhao
	%重庆交通大学硕士学位论文}\fancyhead[CO]{\songti \wuhao  摘要}
%\renewcommand{\headrulewidth}{0.1pt}
\fancyfoot[C]{\songti \xiaowuhao  \thepage}
\renewcommand{\headrulewidth}{0.8pt}

%\addcontentsline{toc}{chapter}{\numberline{中文摘要}}

\pagenumbering{Roman} \topskip 54pt
\centerline{\heiti \sanhao 摘\quad 要} \vskip 1cm

\input{chap/0_abs_CN}

~\\
\noindent{\heiti\xiaosihao 关键词}：{\songti\xiaosihao {***，****，***}}

%%%%%%%%%%%%%%%%%%%%%%%%%%%%%%%%%%%%%%%%%%%%%%%%%%%%%%%%%%%%%%%%%%%%%%%%%%%%%%%%%%%%%%%%%%%%%%%%%%%%%%%%%%%%%%%%%%%%%%%%%%%%%%%%%%%%%%%%%%%%%%%%%%%%%%%%%%%
\newpage
%\addcontentsline{toc}{chapter}{\numberline {英文摘要}}
\centerline{\bf\bf\sanhao ABSTRACT} \vskip 1cm

\input{chap/0_abs_EN}

\quad\\
{\bfseries KEY WORDS:}\quad{***，****，***}

%%%%%%%%%%%%%%%%%%%%%%%%%%%%%%%%%%%%%%%%%%%%%%%%%%%%%%%%%% 绪论
%%%%%%%%%%%%%%%%%%%%%%%%%%%%%%%%%%%%%%%%%%%%%%%%%%%
\newpage
\setlength{\baselineskip}{20pt}

\thispagestyle{empty}

\titlecontents{chapter} [2.5em] {\bfseries \zihao{4}\vspace{-5pt}}
{\contentslabel {2em}}{\hspace* {-4em}}{~\titlerule*[0.4pc] {$\boldsymbol{.}$}~\contentspage}

\titlecontents{section} [4em] {\zihao{5} \vspace{-5pt}}
{\contentslabel {2em}}{\hspace* {-4em}}{~\titlerule*[0.4pc] {$.$}~\contentspage}

\titlecontents{subsection} [6.5em] {\zihao{5} \vspace{-5pt}}
{\contentslabel {2.5em}}{\hspace* {-4em}}{~\titlerule*[0.4pc] {$.$}~\contentspage}


\renewcommand{\contentsname}{目\ \ \ \ 录\vspace{1cm}}
{\setlength{\parskip}{1pt} \tableofcontents
	\thispagestyle{fancy} \fancyhead[CE]{\songti \wuhao
		重庆交通大学硕士学位论文}\fancyhead[CO]{\songti \wuhao  目录}

\newpage
\setcounter{page}{1}
\pagenumbering{arabic}
%%%%%%%%%%%%%%%%%%%%%%%%%%%%%%%%%%%%%%%%%%%%%%%%%%%%%%%%%%%%%%%%%%%%%%%%%%%%%%%%%%%%%%%%%%%%%%%%%%%%%%%%%%%%%%%%%%%%%%%%%%%%%%%%%%%%%%%%%%%%%%%%%%%%%%%%%%%
\newpage
\setlength{\baselineskip}{20pt}

% \chapter{\heiti\qquad 绪论}
\chapter{\heiti{\qquad 绪论}}

\pagestyle{fancy} \fancyhead[CE]{\songti \wuhao
重庆交通大学硕士学位论文}\fancyhead[CO]{\songti \wuhao 绪论}
\indent

\thispagestyle{fancy} \fancyhead[CE]{\songti \wuhao
重庆交通大学硕士学位论文}\fancyhead[CO]{\songti \wuhao 绪论}
\renewcommand{\headrulewidth}{0.1pt}


%%%%%%%%%%%%%%%%%%%%%%%%%%%%%%%%%%%%%

\chapter{\heiti{\qquad   预备知识}}


\fancyhead[CE]{\songti \wuhao
重庆交通大学硕士学位论文}
\fancyhead[CO]{\songti \wuhao 预备知识}
\indent

\thispagestyle{fancy} \fancyhead[CE]{\songti \wuhao
重庆交通大学硕士学位论文}\fancyhead[CO]{\songti \wuhao 预备知识}
\renewcommand{\headrulewidth}{0.1pt}

\input{chap/1_Intro}

%%%%%%%%%%%%%%%%%%%%%%%%%%%%%%%%%%%%%%%%%%%%%%%%%%%%%%%
\chapter{\heiti{\qquad 2}}

\pagestyle{fancy} \fancyhead[CE]{\songti \wuhao
重庆交通大学硕士学位论文}\fancyhead[CO]{\songti \wuhao 2}
\indent

\thispagestyle{fancy} \fancyhead[CE]{\songti \wuhao
重庆交通大学硕士学位论文}\fancyhead[CO]{\songti \wuhao 2}
\renewcommand{\headrulewidth}{0.1pt}
\input{chap/2_content}

%%%%%%%%%%%%%%%%%%%%%%%%%%%%%%%%%%%%%%%%%%%%%%%%%%%%%%%%%

\chapter{\heiti{\qquad 4}}

\pagestyle{fancy}
\fancyhead[CE]{\songti \wuhao 重庆交通大学硕士学位论文}
\fancyhead[CO]{\songti \wuhao 4}
\indent

\thispagestyle{fancy} 
\fancyhead[CE]{\songti \wuhao 重庆交通大学硕士学位论文}
\fancyhead[CO]{\songti \wuhao 4}
\renewcommand{\headrulewidth}{0.1pt}

\input{chap/4_content}


%%%%%%%%%%%%%%%%%%%%%%%%%%%%%%%%%%%%%%%%%%%%%%%%
\chapter{\heiti{\qquad 3}}

\pagestyle{fancy} \fancyhead[CE]{\songti \wuhao
重庆交通大学硕士学位论文}\fancyhead[CO]{\songti \wuhao 3}
\indent

\thispagestyle{fancy} \fancyhead[CE]{\songti \wuhao
重庆交通大学硕士学位论文}\fancyhead[CO]{\songti \wuhao 3}
\renewcommand{\headrulewidth}{0.1pt}


%%%%%%%%%%%%%%%%%%%%%%%%%%%%%%%%%%%%%%%%%%%%%%%%%

\chapter{\heiti{\qquad 总结和讨论}}

\fancyhead[CE]{\songti \wuhao 重庆交通大学硕士学位论文} 
\fancyhead[CO]{\songti \wuhao  总结和讨论}
\indent
\fancyhead[CE]{\songti \wuhao 重庆交通大学硕士学位论文} 
\fancyhead[CO]{\songti \wuhao  总结和讨论}
\renewcommand{\headrulewidth}{0.1pt}
\section{本文总结}

\section{后续工作展望}

%%%%%%%%%%%%%%%%%%%%%%%%%%%%%%%%%%%%%%
\fancyhead[CE]{\songti \wuhao 重庆交通大学硕士学位论文} 
\fancyhead[CO]{\songti \wuhao  参考文献}
\indent
\thispagestyle{fancy} 
\fancyhead[CE]{\songti \wuhao 重庆交通大学硕士学位论文} 
\fancyhead[CO]{\songti \wuhao  参考文献}
\renewcommand{\headrulewidth}{0.1pt}


\printbibliography[heading=bibintoc,title={参考文献\vspace{1cm}}]
	
%%%%%%%%%%%%%%%%%%%%%%%%%%%%%%%%%%%
\thispagestyle{empty}
\ \ \ \

\chapter*{\heiti{致\quad 谢}}
\addcontentsline{toc}{chapter}{\mbox{\heiti \qquad 致谢}}
\pagestyle{fancy} \fancyhead[CE]{\songti \wuhao
    重庆交通大学硕士学位论文}\fancyhead[CO]{\songti \wuhao  致谢}
\indent
\thispagestyle{fancy} \fancyhead[CE]{\songti \wuhao
    重庆交通大学硕士学位论文}\fancyhead[CO]{\songti \wuhao  致谢}

本文是在导师***教授的悉心指导下完成的, 从论文选题, 到撰写完成,
每一步都凝聚着导师的心血, 他的许多重要工作和思想都是本文的基础.
多年以来, *老师以其渊博的知识引领我进入科学研究的前沿,
用敏锐的科学洞察力指引我的研究方向,
用严谨治学的风格和锐意开拓的精神敦促我在学术上的成长,
为我的科研工作打下了坚实的基础. 同时,
*老师在生活上也给了我极大的支持、鼓励和关怀. 在此,
我向*老师表示深深的敬意和衷心的感谢.

我还要衷心感谢我的硕士生导师邹云志教授给予的关心, 鼓励和帮助.

感谢***等所有师兄师姐、师弟师妹们的关心和帮助.

感谢**大学数学学院的领导和老师给予我的关心和支持.
同时也感谢**大学数学学院为我提供的良好的科研条件、宽松的学习环境.

最后，我衷心地感谢在百忙之中审阅我论文和参加答辩的各位专家、教授！
\\

\quad\quad\quad\quad\quad\quad\quad\quad\quad\quad\quad\quad\quad \quad\quad\quad\quad\quad\quad\quad\quad\quad\quad\quad\quad\quad  ~ ******

\quad\quad\quad\quad\quad\quad\quad\quad\quad\quad\quad\quad \quad\quad\quad\quad\quad\quad\quad\quad\quad\quad\quad\quad  2023年5月于重庆

	%%%%%%%%%%%%%%%%%%%%%%%%%%%%%%%%%%%%%%%%%%%%%%%%%%%%%%%%%%%%%%%%%%%%%%%%%%%%%%%%%%%%%%%%%%%%%%%%%%%%%%%%%%%%%%%%%%%%%%%%%%%%%%%%%%%%%%%%%%%%%%%%%%%%%%%%%%%%%%%%%%%%%%%%%%%
	%%%%%%%%%%%%%%%%%%%%%%%%%%%%%%%%%%%%%%%%%%%%% 声明
	%%%%%%%%%%%%%%%%%%%%%%%%%%%%%%%%%%%%%%%%%%%%%
	
	
\thispagestyle{empty}
\ \ \ \


\topskip 80pt
\chapter*{\heiti{攻读硕士学位期间发表的学术论文及参与的科研项目情况}}
\addcontentsline{toc}{chapter}{\mbox{\heiti \qquad 攻读硕士学位期间发表的学术论文及参与的科研项目情况}}

\thispagestyle{fancy} \fancyhead[CE]{\songti \wuhao
    重庆交通大学硕士学位论文}\fancyhead[CO]{\songti \wuhao  攻读硕士学位期间发表的学术论文及参与的科研项目情况}

\indent

\section*{1\ 发表的学术论文}
\addcontentsline{toc}{section}{\mbox{ \qquad 1\ 发表的学术论文}}

%1. Robust stability of fractional-order quaternion-valued neural networks with neutral delays and parameter uncertainties[J]. Neurocomputing, 2021, 420: 70-81. (第二作者，SCI 二区，已检索)
%
%2. Global Mittag-Leffler stability for fractional-order quaternion-valued neural networks with piecewise constant arguments and impulses[J]. International Journal of Systems Science. DOI: 10.1080/00207721.2021.2023688. (第一作者，SCI三区，已检索)

\noindent [1] \textbf{Li X}, Song Q, Zhao Z, Liu Y, Alsaadi F E, Optimal control and zero-sum differential game for Hurwicz model considering singular systems with multifactor and uncertainty [J], International Journal of Systems Science, 2021, 53(7), 1416-1435. (对应论文第六章内容，SCI 检索, 已发表)

\vspace{0.2cm}

\noindent [2] \textbf{Li X}, Song Q, Liu Y, Alsaadi F E, Nash equilibrium and bang-bang property for the non-zero-sum differential game of multi-player uncertain systems with Hurwicz criterion [J], International Journal of Systems Science, 2022, 53(10), 2207-2218. (对应论文第三章内容，SCI 检索, 已发表)

\vspace{0.2cm}





\section*{2\ 参与的科研项目}
\addcontentsline{toc}{section}{\mbox{ \qquad 2\ 参与的科研项目}}

\noindent [1] 国家自然科学基金，62176032，四元数斜域上的伊藤公式及其在神经网络稳定性分析中的应用，2022年1月-2025年12月，在研，参与。

\end{document}
